\chapter{Resumen}


\textbf{Resumen:} 

El huevo es uno de los alimentos principales de las personas, con un consumo per cápita de 22 kilos anuales, en México, siendo el 4° productor mundial con más de 2.77 millones de toneladas anuales \cite{AGRICULTURA}. Esto ha provocado que su proceso sea cada vez más automatizado. Sin embargo, también existen los productores para consumo personal, que lo hacen montando un sistema rural o de traspatio. El presente documento propone y diseña un sistema automático para el monitoreo de la alimentación y la producción de gallina de postura, para un sistema de producción de traspatio, que pueda apoyar al bienestar animal de estas y la eficiencia de la producción. Así mismo, se desarrolla el diseño con base en los criterios planteados para implementación empleando la metodología de diseño adecuada. Presentando y analizando los diferentes antecedentes que pueden existir a un objetivo en común al que se propone en el proyecto, se presenta la solución de este. Se busca que el sistema no sea invasivo y que pueda mejorar la producción, reduciendo gastos.



\textbf{Palabras Clave:} Gallina de postura, monitoreo, alimentador animales, RFID.



%%%%%%fin del archivo

\endinput 